\documentclass[../../main.tex]{subfiles}
\begin{document}
\begin{oframed}
  \subfile{problem}
\end{oframed}

{\bf [解答]}

\paragraph{(1)}

与えられた漸化式より，
\begin{align}
  &a_n<a_{n+1}\nonumber\\
  \Longleftrightarrow&2a_n<a_n^2+1\nonumber\\
  \Longleftrightarrow&(a_n-1)^2>0\nonumber\\
  \Longleftrightarrow&a_n\not=1\label{eq:1}
\end{align}
である．$x\not=2$ならば，$a_1\not=1$だから，漸化式より帰納的に
$a_n\not=1$である．これと\cref{eq:1}より題意は示された．\qed

\paragraph{(2)}

漸化式および$x>0$から，$a_n>0$であることに注意すると，$a_{n+1}$は$a_n$の単調増加関数である．
$x<2$ならば$a_1<1$である．そこで以下$k\in\mathbb{N}$に対して$a_k<1$が成立すると仮定すると
\begin{align}
  a_{k+1}=\dfrac{a_k^2+1}{2}<\dfrac{1+1}{2}=1
\end{align}
故に帰納法により$\forall n,a_n<1$となる．\qed

次に，後半部分を考える．題意より
\begin{align}
  \left\{
    \begin{array}{ll}
      1-\dfrac{x}{2}>\epsilon\\
      a_k\le 1-\epsilon&k=1,2,\cdots,n
    \end{array}\right.\label{eq:2}
  \end{align}
  と仮定する．第一式から$a_1<1-\epsilon$が保証される．

  ここで漸化式より，$k\le n$のとき
  \begin{align}
    &2(a_{k+1}-a_k)=(a_k-1)^2 \ge \epsilon^2 \quad (\because\cref{eq:2})
  \end{align}
  となる．ただし$a_k<1$より$a_k-1 \le -\epsilon < 0$であることを用いた．
  $k=1,2,\cdots,n$について和をとって，
  \begin{align}
    & 2(a_{n+1}-a_1)=\sum_{k=1}^n (a_k-1)^2 \ge n\epsilon^2 \\
    \therefore
    & a_{n+1}-\dfrac{x}{2} \ge \dfrac{n\epsilon^2}{2}
  \end{align}
  前半部分から$a_{n+1}<1$だから，
  \begin{align}
    & 1-\dfrac{x}{2}>\dfrac{n\epsilon^2}{2} \\
    \therefore
    & 2-x>n\epsilon^2
  \end{align}
  となって，題意の不等式を得る．\qed
  \newpage
  \end{document}
